\chapter{Discussion and Conclusion}
\label{chap:conclusion}
\setlength{\emergencystretch}{1.5em}

Chapter~\ref{chap:results} reported what the experimental evidence
shows. This chapter interprets those findings in relation to the three
research hypotheses introduced in
Chapter~\ref{chap:problem}, including how the evidence supports a more
scenario-dependent interpretation than originally hypothesised. It then
discusses the broader implications of the proposed framework, its
limitations, and directions for future research.

% ============================================================
\section{Summary of the Work}
\label{sec:conclusion_summary}
% ============================================================
\begin{sloppypar}
This thesis investigated whether a unified reinforcement learning
agent, jointly optimising energy dispatch and diesel inventory
ordering, provides measurable operational advantages over decoupled
classical and optimisation-based approaches, and how those advantages
depend on delivery-logistics conditions. The gap motivating this work is
practical: prior RL-based energy management for remote base stations
routinely models battery dispatch and generator scheduling but omits
the upstream replenishment decision --- when and how much diesel to
order --- despite delayed replenishment being an important source of
reliability risk in fuel-constrained deployments.
\end{sloppypar}

The problem was framed as a constrained Markov decision process in
which a single policy (RLInv) controls generator commitment and diesel
order quantity jointly, with battery dispatch handled by a fixed
priority rule and an 11-dimensional observation vector including
explicit inventory and pipeline state (Chapter~\ref{chap:problem}).
RLInv was evaluated against eight comparator policies spanning
rule-based, reinforcement learning, and optimisation-based approaches ---
including, notably, a receding-horizon MPC baseline and an
Oracle-MPC optimistic dispatch reference under perfect future load, solar-generation and grid-availability information, added specifically to distinguish whether any
observed advantage reflects better \emph{ordering} or merely better
\emph{dispatch} (Chapter~\ref{chap:methodology}). The final campaign
evaluated all nine core policies across 10 ITU/Zindi sites and four
delivery-logistics scenarios of increasing severity (Normal, Delayed,
Monsoon, Extreme). The five learned policies (RLInv, TrackB, A5, A6,
A7; see Table~\ref{tab:ablation-design}) used 10 independent training
seeds, with all policies evaluated
under the corresponding matched protocol and paired training-seed
statistical comparisons for the primary policy and ablation analyses
(Chapter~\ref{chap:experimental-execution}).

% ============================================================
\section{Key Results and Discussion}
\label{sec:conclusion_results}
% ============================================================

\subsection{\revised{What the combined controller comparisons show}}

\begin{revisedblock}
Within the original comparison set, RLInv achieved the lowest
mean Expected Energy Not Served (EENS) in all four logistics
scenarios. The expanded comparisons qualify that result.
Calibrated B1*-HighBuffer achieved lower mean EENS under
Delayed and Extreme logistics, with a nearly identical mean
under Monsoon. The supplementary always-order-max controller,
which requests diesel generation whenever feasible and places
the largest feasible discrete replenishment order, matched or
improved on RLInv's mean EENS in all four scenarios while
placing slightly fewer accepted orders. This comparison does
not establish lower procurement cost because the number of
accepted orders does not measure ordered fuel volume,
holding cost or expenditure.

The simulator's fixed PV--grid--battery dispatch priority
leaves a site-specific residual demand, which is partitioned
between diesel energy and unmet load
(Equation~\ref{eq:null_fixed_residual}). On the ten-site
360-hour evaluation, no diesel dispatch gives mean EENS of
222.932\,kWh; greedy diesel dispatch without replenishment
reduces it to 86.370\,kWh. Allowing maximum feasible
replenishment reduces it further to 0.019, 2.517, 9.891
and 51.745\,kWh under Normal, Delayed, Monsoon and Extreme
logistics, respectively. The 136.562\,kWh dispatch contrast
is identical across scenarios because the no-order reference
controllers do not depend on supplier lead time.

The implemented energy reward can be written, using
$U_t=R_t-D_t$, as
$(99D_t-100R_t-0.4G_t)/\bar d$, in addition to stockout
and physical-violation penalties. This gives a strong
immediate incentive to supply diesel energy when it avoids
unmet demand. As no explicit order-placement, ordered-volume
or holding costs are included, a simple aggressive
replenishment rule is a strong empirical comparator. The
identity and reward expression explain the narrow reliability
gaps observed among strong controllers, but do not prove the
rule globally optimal or uniquely identify why RLInv did not
match it.

Under the supplementary site-normalised reward comparison
reported in Chapter~\ref{chap:results}, RLInv captures
approximately 100.0\%, 99.5\%, 99.4\% and 98.9\% of the
observed improvement from the no-DG/no-order reference to
the highest-scoring evaluated controller across the four
scenarios. These descriptive percentages show a small
remaining gap relative to the observed reference-to-best
range, not a theoretical bound on achievable performance.
\end{revisedblock}


\subsection{Why the mechanism results are scenario-dependent, not uniform}

The three pre-specified hypotheses were tested through paired,
seed-level ablations rather than a single aggregate comparison, and the
honest answer in each case is more nuanced than a simple confirmed/rejected
verdict.

\begin{revisedblock}
\textbf{Joint replenishment and dispatch learning (H1).}
The original RLInv--A6 ablation found lower RLInv EENS under
Normal and Extreme logistics, but not a statistically
distinguishable advantage under Delayed or Monsoon conditions.
This is evidence for a configuration-dependent benefit of the
unified learned controller over the tested decomposed learned
dispatch configuration; the two also differ in dispatch action
space. The later always-order-max result sharpens the
interpretation: a non-learning rule jointly determines diesel
dispatch and replenishment and matches or improves on RLInv's
mean EENS. Thus, coordinating dispatch with replenishment is
important in this environment, but the experiments do not
establish that learning that coordination provides additional
reliability over a strong feasible rule.
\end{revisedblock}

\textbf{Action masking (H2).} A6 substantially outperforms TrackB
under Monsoon and Extreme logistics. However, the comparison differs
in both masking and observation design (Chapter~\ref{chap:methodology}), so the gap cannot be
attributed to masking alone. The cleaner RLInv--A7 comparison, which
retains the same observation and joint-control architecture, finds no
statistically distinguishable mean-EENS effect from removing masking.
The evidence therefore does not establish a universal independent
performance benefit from masking, although masking remains the
mechanism used to enforce action feasibility. The observed A6--TrackB
gap is architecture-dependent and appears only under stressed
logistics.

\textbf{Inventory observability (H3).} RLInv achieves numerically lower
EENS than an ablation with inventory and ordering-history observations
removed, in every scenario, but the difference never reaches
significance (closest at Extreme, $p=0.058$). The result is interpreted
cautiously. This may indicate that the remaining observations (battery
state-of-charge, load, solar generation, grid availability, and
time-of-day features) provide useful indirect information, or that the
incremental contribution of the five removed channels is small under
the tested conditions. Because the five removed logistics-related
channels are removed together, the experiment does not isolate the
contribution of inventory level from pending-order and
ordering-history information.

\subsection{Behaviour, robustness, and the exploratory extension}

Behavioural analysis shows RLInv's ordering frequency decreasing
monotonically from 2.16 orders per episode under Normal logistics to
1.34 under Extreme, consistent with the policy reducing new orders as
delivery completion within the evaluation horizon becomes increasingly
unlikely, while stockout events remain at zero even under Extreme
stress --- non-zero EENS occurs without recorded diesel stockout or
hard constraint violation. A sensitivity analysis across a wider
sequence of mean delivery
lead times shows RLInv's advantage over the best classical baseline
peaking at 29.5\,kWh around 72\,h and narrowing to 8.6\,kWh by 336\,h;
this narrowing reflects a physical constraint common to every policy
(tank capacity bounds how much can be pre-ordered regardless of
strategy), not RLInv ceasing to function.

RLInv's performance is robust to the tested reward-coefficient choices
(EENS varies by only 3.32\,kWh across the evaluated $\lambda$ range of
25--500), to tank capacity (RLInv's advantage narrows as tank capacity
increases, showing that sufficiently large physical buffers reduce the
value of proactive learned ordering), and to the parametric form
assumed for delivery delay (policy rankings are preserved under
lognormal and Weibull alternatives, with RLInv less sensitive to this
modelling choice than the classical baseline). One genuine limitation
was found under stochastic (rather than deterministic) grid outages: a
policy trained under this harder condition improves under Normal
logistics but degrades under Delayed, Monsoon, and Extreme, most
severely at Monsoon (+43.4\%). This pattern suggests a possible
trade-off between adaptation to grid uncertainty and fuel conservation
under prolonged delivery
disruption, reported as an identified limitation rather than resolved.

An exploratory extension (RLInv-DA) investigated whether explicit
delay-history tracking, motivated by Phase-I reviewer feedback,
improves on RLInv further. Despite a verified, working mechanism (a
correctly updating delay estimate under a persistent two-regime
supplier model), the extension produced no consistent EENS improvement
over RLInv on the narrower evaluation grid used for this exploratory
study. This may indicate that the existing state already
contains sufficient indirect information for the tested conditions, or
that the simple EWMA signal does not capture the relevant supplier
dynamics; the added complexity is not justified at the scale evaluated.

% ============================================================
\section{Thesis Contributions}
\label{sec:conclusion_contributions}
% ============================================================

This thesis makes the following contributions to RL-based energy
management of remote telecom infrastructure:

\begin{enumerate}

  \item \textbf{A CMDP formulation with explicit, jointly controlled
  inventory and dispatch decisions.} The inclusion of diesel inventory
  and in-transit pipeline state alongside battery state-of-charge and
  grid availability enables anticipatory ordering that reacts to
  lead-time uncertainty before a stockout occurs, and is evaluated
  against eight rule-based, RL-based and optimisation-based comparator
  policies.

  \item \textbf{An empirical comparison of coupled and decomposed
  control.} The study compares RLInv with rule-based, RL-based, and
  optimisation-based alternatives. The A6 ablation provides a cleaner
  practical assessment of learned replenishment, while its reduced
  dispatch action space is acknowledged as a remaining design
  difference. MPC and Oracle-MPC help contextualise dispatch and
  forecasting quality; they do not by themselves isolate ordering. The
  result is a genuinely
  scenario-dependent finding, reported honestly rather than as a single
  uniform verdict.

  \item \textbf{A qualified assessment of action masking.} Action
  masking is implemented as deterministic feasibility filtering. The
  cleaner RLInv--A7 comparison finds no statistically distinguishable
  mean-EENS effect when ordering is jointly learned, while the larger
  A6--TrackB gap under stressed logistics is confounded by differences
  in observation design. This distinction prevents feasibility
  enforcement from being mistaken for a universal performance gain.

  \item \textbf{A reproducible multi-stage evaluation framework.} The
  primary comparison and core ablations use the common 10-site,
  10-seed, and four-scenario protocol. Additional reward, tank,
  distributional, stochastic-grid, and delay-aware studies use
  explicitly documented experimental grids appropriate to their
  respective questions. Canonical source files, aggregation rules and
  statistical units are documented for each analysis to support
  reproducibility.

  \item \textbf{An exploratory delay-aware extension.} The
  delay-aware extension is reported as a verified exploratory study
  that did not produce a consistent improvement under the tested
  conditions.

\end{enumerate}

% ============================================================
\section{Critique and Limitations}
\label{sec:conclusion_limitations}
% ============================================================

\revised{The following limitations define the scope of the
experimental findings.}
\begin{sloppypar}
\begin{description}

\item[\revised{Conservation-linked reliability measures.}]
\revised{For the deterministic-grid evaluation with fixed
PV--grid--battery dispatch priority, the site-specific residual
satisfies $\mathrm{EENS}_i+E_{\mathrm{diesel},i}=R_i$ to
numerical precision. EENS and diesel generation therefore do
not represent independent reliability and efficiency gains.
This is a restriction on the measured outcomes, not proof that
the sequential inventory-control problem has only one action
or one decision variable.}

\item[\revised{Reward specification and replenishment economics.}]
\revised{The implemented reward strongly prioritises avoided
unserved energy; after residual substitution, its energy
component is $(99D_t-100R_t-0.4G_t)/\bar d$.
It does not explicitly charge for order placement, ordered
fuel volume or inventory holding. The comparison therefore
cannot establish an economic optimum or a general fuel-cost
advantage. The simplified reward makes frequent feasible
replenishment a strong baseline, but does not analytically
establish a globally optimal ordering policy.}

\item[\revised{Limited separation among strong controllers.}]
\revised{Nineteen of the forty site--scenario combinations in
the four-algorithm comparison produced exact EENS ties.
The supplementary reward comparison places RLInv close to
the strongest observed controller relative to the no-DG
reference. These results constrain the discrimination
possible with the current evaluation metric and fixed dispatch
priority; they do not show that the reward is flat or that all
controllers behave equivalently.}

\item[\revised{Post-hoc null-controller coverage.}]
\revised{The original planned comparison omitted a simple
no-generator/no-replenishment floor and an aggressive feasible
replenishment rule. These controllers were added subsequently
as a post-hoc diagnostic, not as preregistered primary
comparators. Their evaluation seed blocks are not independent
training replicates, so the new comparisons are interpreted
descriptively rather than as confirmatory superiority tests.}

\item[\revised{Restricted dispatch and replenishment actions.}]
\revised{PV, grid and battery dispatch follow fixed priority
rules; RLInv controls generator commitment and a six-action
combination of generator state and tank-relative order size.
The order classes are 0\%, 30\% and 60\% of tank capacity,
with at most one order in transit. The experiments therefore
do not establish performance under co-optimised battery
scheduling, continuous ordering or full economic dispatch.}

\item[\revised{Synthetic supplier logistics.}]
\revised{The ITU/Zindi traces provide site energy inputs but
not observed diesel order or delivery records. Lead-time
models and pipeline states are simulation overlays; even the
lognormal and Weibull robustness checks compare assumed
processes rather than real supplier data.}

\item[\revised{Algorithm-comparison and training differences.}]
\revised{MaskablePPO, unmasked PPO, DQN and A2C were compared
using five training seeds, but the training initial-inventory
settings were not identical across all methods. Many site
conditions were non-discriminating, and no paired difference
survived multiplicity correction. The results do not establish
superiority of an RL algorithm family.}

\item[\revised{Transfer and deployment scope.}]
\revised{Zero-shot transfer used three independently trained
shared policies and five held-out sites drawn from the same
60-day dataset. Approximately fivefold fewer nominal training
interactions do not establish a comparable reduction in
wall-clock time or real deployment cost. Every tower was
simulated independently; shared suppliers, transport capacity,
fleet coordination and out-of-region or seasonal transfer
were not evaluated.}

\item[\revised{Narrower exploratory robustness studies.}]
\revised{The delay-aware RLInv-DA extension used only three
sites and three training seeds and produced no consistent
benefit. Stochastic-grid training revealed scenario-dependent
trade-offs that remain unresolved. These supplementary results
should not be treated as if they had the full original
experimental population.}

\end{description}
\end{sloppypar}

% ============================================================
\section{Future Work}
\label{sec:conclusion_future}
% ============================================================

The remaining research directions follow directly from the
limitations and conditional findings identified in this thesis. Four
major directions are discussed below, followed by two smaller
extensions.

\textbf{1.\ Richer logistics and supplier-behaviour modelling.} Two
results in this thesis point the same way. First, RLInv-DA showed that
a simple online estimator (EWMA) of realised delivery delay does not
improve replenishment at the evaluation scale tested. Future work could
compare richer state-estimation approaches, including Bayesian filters,
state-space models, and learned delay predictors, using empirical
supplier data where available.
Second, training under stochastic grid outages improves Normal-scenario
performance at the cost of degraded performance under logistics stress;
a natural extension is to train jointly under stochastic grid
\emph{and} varied logistics scenarios simultaneously, testing whether
this trade-off is an artefact of modelling only one source of
uncertainty at a time. Both directions aim at the same underlying gap:
richer, more realistic supplier and infrastructure behaviour than the
geometric and two-state abstractions used here.

\textbf{2.\ Fleet-level coordination.} Realistic
deployments often involve several towers sharing a common diesel
supplier and competing for delivery capacity during regional
disruptions. The single-tower CMDP formulated here cannot capture this
coordination problem. Future work could model shared supplier capacity
through multi-agent reinforcement learning, centralised fleet
optimisation, or hybrid coordination approaches, as the natural
fleet-level counterpart to the single-tower joint-optimisation result
established here.

\textbf{3.\ Validation against empirical delivery datasets.} Every
distributional-robustness check in this thesis (geometric, lognormal,
Weibull) compares parametric modelling assumptions against one another,
not against observed delivery logs, since none exist in the available
dataset. Partnering with a telecom operator or fuel logistics provider
to obtain real delivery-time records would allow the lead-time model
itself to be validated empirically, rather than only stress-tested
across plausible parametric forms.

\begin{sloppypar}
\textbf{4.\ Transfer learning and cross-region deployment.} The
per-site evaluation across multiple profiles in this thesis trains and
evaluates a separate
policy per site; it is a per-site robustness check, not a test of
whether a policy trained on a subset of sites transfers to unseen ones
or to a different deployment region. Training on a subset of sites and
evaluating on held-out sites, or on sites from a different geographic
region with different solar, load, and outage characteristics, would
test this directly and is the natural precursor to any claim of
practical, low-cost deployment across a large or geographically diverse
fleet.

Three smaller extensions follow directly from the objective function and
state representation already defined in this thesis, which could be
explored within the existing simulator framework after adding the
required parameters or forecast inputs. 
\end{sloppypar}
\textbf{Carbon and multi-objective
reporting}: future work could extend the current energy- and
reliability-based objective with explicit emissions and
battery-degradation terms, and report EENS--cost--CO$_2$ Pareto
trade-offs rather than relying on a single scalar reward. \textbf{Solar forecast integration}: RLInv
currently observes only realised PV output; augmenting the observation
with a short-horizon solar forecast, available in practice from
numerical weather prediction services, would let the agent anticipate
low-PV periods rather than reacting to them, a natural complement to
the delay-aware extension already explored on the ordering side.
\textbf{Forecast-aware control.} Forecast-aware dispatch was also
explored during the course of this research. Preliminary experiments
incorporated short-term load and solar forecasts into the MPC
baseline's dispatch planning, to test whether predictive information
about future conditions would improve performance relative to
persistence-based planning. Although this direction showed promise
conceptually, it did not consistently improve operational performance
within the scope of this study. This finding is consistent with RLInv's
design choice of learning directly from the current system state rather
than depending on external forecasts, though forecast-integrated
reinforcement learning remains an interesting direction for future work.

% ============================================================
\section{Closing Remarks}
\label{sec:conclusion_closing}
% ============================================================

\begin{revisedblock}
This thesis examined whether learning diesel replenishment
and generator dispatch jointly improves service reliability
relative to decomposed and classical control. The original
comparisons showed advantages over the initial baseline
configurations, but the expanded tests provide a narrower
conclusion: calibrated classical policies and a simple
feasible-replenishment rule match or improve on RLInv in
several evaluated conditions. Joint coordination of
replenishment and dispatch is operationally relevant;
the experiments do not establish that learning that
coordination yields additional reliability over the
strongest evaluated simple controller.

The explanation lies partly in the implemented benchmark.
With PV, grid and battery dispatch following fixed
priority rules, remaining demand is partitioned between
diesel generation and unserved energy. The reliability-heavy
reward strongly favours available diesel, without explicit
ordering, volume or holding charges. In this setting,
RLInv captured nearly all of the observed improvement over
the no-DG reference, but a simple aggressive replenishment
rule achieved a slightly better observed result. The
contribution is therefore an inventory-coupled sequential
control formulation, an auditable comparative benchmark,
and a quantitative account of its discriminating limits,
not a claim that reinforcement learning is universally
superior. Determining when learning adds economic value
requires richer control decisions, explicit replenishment
costs and validation against realistic supplier constraints.
\end{revisedblock}

Energy security at remote telecom towers is not a niche engineering
problem. In many low- and middle-income countries, reliable
base-station operation supports access to emergency communication,
mobile banking, and public health services, and diesel logistics
account for a substantial share of operating cost in weak-grid
deployments. A controller that is
reliable under routine conditions and whose limits under logistics
stress are precisely characterised, rather than assumed away, is a
more trustworthy basis for real deployment decisions than one whose
advantages are reported without qualification.